% =============================================================
%  preamble.tex —— 统计力学自学讲义 · 导言区
%  编译方式：xelatex ×2（见 build.sh）
% =============================================================

% ---------------- 页面版式 ----------------
\usepackage{geometry}
\geometry{a4paper,left=2.6cm,right=2.6cm,top=2.8cm,bottom=2.8cm}

% ---------------- 数学 ----------------
\usepackage{amsmath}
\usepackage{amssymb}
\usepackage{bm}
\usepackage{mathrsfs}

% ---------------- 图、表 ----------------
\usepackage{graphicx}
\usepackage{booktabs}
\usepackage{array}
\usepackage{tabularx}
\usepackage{xltabular}   % 可跨页的 tabularx（附录 A、B 的长表用）
\usepackage{multirow}

% ---------------- 其他 ----------------
\usepackage{enumitem}
\usepackage{xcolor}
\usepackage{fancyhdr}
\usepackage[unicode]{hyperref}
\hypersetup{
  colorlinks = true,
  linkcolor  = blue!50!black,
  urlcolor   = blue!50!black,
  bookmarksnumbered = true
}

% ---------------- 编号：公式、图、表均按章编号 ----------------
% 与原文一一对应：公式 (1.45)、图 1.1、表 1.1
\numberwithin{equation}{chapter}

% 章号用阿拉伯数字（ctex 默认中文数字，会使公式编号变成「(一.45)」）
\ctexset{chapter/number = \arabic{chapter}}

% ---------------- 双语标题 ----------------
% 正文标题：中文一行、英文另起一行（字号略小）
% 目录条目：中文（English），一行
\newcommand{\bichapter}[2]{%
  \chapter[#1（#2）]{\texorpdfstring{#1\\[0.5em]{\large\mdseries #2}}{#1（#2）}}%
  \markboth{#1}{}%
}
\newcommand{\bisection}[2]{%
  \section[#1（#2）]{\texorpdfstring{#1\\[0.25em]{\normalsize\mdseries #2}}{#1（#2）}}%
}
\newcommand{\bisubsection}[2]{%
  \subsection[#1（#2）]{\texorpdfstring{#1\\[0.2em]{\small\mdseries #2}}{#1（#2）}}%
}

% ---------------- 术语 ----------------
% 原文中意大利斜体的术语：首次出现处写成「中文（English）」，
% 英文保持斜体，并埋下锚点供附录 B 用 \pageref 生成页码。
% 用法：\term{广延量}{extensive}{1.extensive}
\newcommand{\term}[3]{#1（\textit{#2}）\label{term:#3}}

% ---------------- 附录 B 英文列字体 ----------------
% 附录 B 的英文术语统一用 Times New Roman 正体（不用意大利斜体）。
% 生成脚本 _extract/make_appendixB.sh 把每条英文术语包在 {\glossaryfont ...} 中。
\newfontfamily\glossaryfont{Times New Roman}

% ---------------- 非精确微分记号 ----------------
% 原文中热量/功的微小量写作 đQ、đW（d 上加一横杠）。此处先试直接输入该字符；
% 若所用字体缺字，则退回「d 上加斜杠」的画法。
\newcommand{\inexd}{\text{đ}}
% \newcommand{\inexd}{\not{\mathrm{d}}}   % 备用：字体缺 đ 时启用

% ---------------- 页眉页脚 ----------------
\pagestyle{fancy}
\fancyhf{}
\fancyhead[L]{\small\nouppercase{\leftmark}}
\fancyhead[R]{\small\thepage}
\renewcommand{\headrulewidth}{0.4pt}
\fancypagestyle{plain}{%
  \fancyhf{}%
  \fancyfoot[C]{\small\thepage}%
  \renewcommand{\headrulewidth}{0pt}%
}

% ---------------- 图注辅助 ----------------
% 图片由用户自行截取后按 fig/README.md 的命名放入 fig/ 目录（fig<章>-<图>.png/jpg/pdf）。
% 文件存在则自动插入图片；尚不存在时只显示占位框，保证任何时刻都能编译通过。
\newcommand{\figfile}[1]{%
  \IfFileExists{fig/#1.png}{\includegraphics[width=0.8\textwidth]{fig/#1.png}}{%
  \IfFileExists{fig/#1.jpg}{\includegraphics[width=0.8\textwidth]{fig/#1.jpg}}{%
  \IfFileExists{fig/#1.pdf}{\includegraphics[width=0.8\textwidth]{fig/#1.pdf}}{%
  \fbox{\parbox[c][3.6cm][c]{0.8\textwidth}{\centering\small [ 图片占位：#1 ]}}}}}}
% 用法：\cnfigure{fig1-1}{中文图注}
\newcommand{\cnfigure}[2]{%
  \begin{figure}[htbp]
    \centering
    \figfile{#1}%
    \caption{#2}%
    \label{fig:#1}%
  \end{figure}%
}

% ---------------- 卷首摘要（内容留空，待用户最后撰写）----------------
\newenvironment{cnabstract}{%
  \clearpage
  \chapter*{\texorpdfstring{摘要\\[0.5em]{\large\mdseries Abstract}}{摘要（Abstract）}}%
  \phantomsection
  \addcontentsline{toc}{chapter}{摘要（Abstract）}%
  \markboth{摘要}{}%
  \vspace{1em}%
}{%
  \clearpage
}
